\subsection{I-34. Prime cocharacters with a bounded local nullity}
\label{I-34}

\paragraph{Statement and scope.}
Let $n\ge1$, let $\delta\ge1$ be an integer, and consider
\[
 f(\mathbf X)=\sum_{j=1}^{w}\prod_{i=1}^{n}F_{ij}(X_i),
 \qquad 0\ne F_{ij}\in F[X_i],\quad \deg F_{ij}\le\delta.
 \tag{I.67}
\]
Zero product terms can first be omitted. Scalars and signs can be
absorbed into the factors. Let $A_i$ be the $(\delta+1)\times w$
matrix whose $j$th column is the coefficient vector of $F_{ij}$.
We call $w-\operatorname{rank}A_i$ its local nullity.

There is an explicit family of polynomial-magnitude positive
integer cocharacters with the following properties, over every
characteristic.

If some $A_i$ has local nullity at most $c$, the family changing
at most $k=\min\{n,\max\{1,c\}\}$ coordinates of a base prime
weight vector hits every nonzero $f$ in (I.67). For fixed $c$
this gives polynomial sampling size even when $w$ grows. The
promise itself requires $w-c\le\delta+1$; the assertion does not
silently treat $c$ as constant for arbitrary representations.

Without a local-nullity promise, sums of at most three separated
products require only one changed coordinate. More generally,
sums of $w\ge3$ products are hit using at most
$\min\{n,w-2\}$ changed coordinates. The last bound has an
exponent depending on $w$ and is not a polynomial-in-width PIT
result. Constant-top-fan-in PIT is already known, for example
from I-B08. The purpose here is to test and repair the specific
prime-cocharacter construction, not to claim a new general
complexity bound for an established constant-width class.

One-coordinate variation cannot extend this result to all sums
of four products: an explicit tensor-rank-four multilinear
example below vanishes on every entire one-coordinate cylinder
around an admissible prime curve. Thus increasing the number of
alternate primes while still varying only one coordinate does
not remove that obstruction.

\paragraph{1. The explicit family and its sampling cost.}
Set
\[
 H=2\delta(n-1),\qquad
 M=H+n(\delta+1)+2,\qquad B=16M^2.
\]
Choose $n(\delta+1)$ distinct primes in $(H,B]$, organized as
$p_i$ and $q_{i,a}$ for $i\in[n]$, $a\in[\delta]$.
The elementary prime count in I-33 applies verbatim: it gives
$\pi(B)\ge2M$, and at most $H$ primes lie below the threshold.
All these primes can be found by a polynomial-time sieve.

Let $\mathcal W_k$ contain the base vector $(p_1,\ldots,p_n)$
and every vector obtained by replacing $p_i$ by one of
$q_{i,1},\ldots,q_{i,\delta}$ on at most $k$ coordinates.
Then
\[
 |\mathcal W_k|=
 \sum_{a=0}^{k}\binom na\delta^a.
 \tag{I.68}
\]
Each curve $X_i=T^{\omega_i}$ has substituted degree at most
$n\delta B$. A scalar grid of $n\delta B+1$ distinct values
therefore gives a hitting set of size at most
\[
 (n\delta B+1)\sum_{a=0}^{k}\binom na\delta^a.
 \tag{I.69}
\]
An adequate working extension is understood if the coefficient
field is too small. These are degree bounds in the formal
indeterminate $T$, so no characteristic restriction on the
integer weights is needed.

The reason to use $\delta$ alternate primes, rather than just
one, is exact interpolation over $F(T)$. For $S\subseteq[n]$
write
\[
 f_S(\mathbf Z_S,T)=
 f\bigl((Z_i)_{i\in S},(T^{p_i})_{i\notin S}\bigr).
 \tag{I.70}
\]
This is the coordinate cylinder with the coordinates in $S$
free. It has individual degree at most $\delta$ in each $Z_i$.
The $\delta+1$ elements
$T^{p_i},T^{q_{i,1}},\ldots,T^{q_{i,\delta}}$ are distinct
in $F(T)$. Consequently, if every curve from $\mathcal W_k$
kills $f$, tensor-product interpolation over $F(T)$ shows that
$f_S=0$ for every $|S|\le k$. Conversely, if all these cylinders
are zero, every such curve is zero. This equivalence uses only
formal roots and is valid in all characteristics.

\paragraph{2. Nullity one: the circuit argument.}
We prove that if one $A_i$ has nullity at most one, then a
nonzero $f$ has a nonzero one-coordinate cylinder. Suppose
instead that every such cylinder is zero. Put
\[
 B_j(T)=\prod_{\ell\ne i}F_{\ell j}(T^{p_\ell}).
\]
Every $B_j$ is nonzero. The identity $f_{\{i\}}=0$ says
$A_i(B_1,\ldots,B_w)^T=0$ over $F(T)$.
If $A_i$ has full column rank this is impossible. Otherwise
its kernel is one-dimensional and defined over $F$, so
\[
 B_j(T)=h(T)c_j\quad(1\le j\le w),
 \qquad h\ne0,\quad c_j\in F^\times.
 \tag{I.71}
\]
The nonvanishing of every $c_j$ follows from that of every
$B_j$. Thus the columns of $A_i$ form a circuit: their unique
relation, up to a scalar, has full support.

For each $j$, the ratio $B_j/B_1$ is constant. Apply the
divisor separation lemma I-33 to the $n-1$ remaining variables.
Its threshold is no larger than our $H$; empty products when
$n=1$ cause no problem. We obtain
\[
 \frac{F_{\ell j}(X)}{F_{\ell1}(X)}
       =d_{\ell j}X^{b_{\ell j}}
 \quad(\ell\ne i),
 \qquad d_{\ell j}\in F^\times,
 \quad b_{\ell j}\in\mathbb Z.
 \tag{I.72}
\]
The constants lie in $F$: the ratio has coefficients in $F$
and equals a constant times a monomial. Negative exponents
are allowed in this rational-function identity. Set
$b_{\ell1}=0$.

Fix $\ell\ne i$ and use the zero cylinder with $X_\ell=Z$.
Dividing by nonzero rational functions and using (I.71)--(I.72)
reduces it to the Laurent identity
\[
 \sum_{j=1}^w c_j F_{ij}(T^{p_i})
       (Z/T^{p_\ell})^{b_{\ell j}}=0.
 \tag{I.73}
\]
For every nonempty group $J_b=\{j:b_{\ell j}=b\}$, its
coefficient gives
$\sum_{j\in J_b}c_jF_{ij}(T^{p_i})=0$.
The substitution $X\mapsto T^{p_i}$ is injective on $F[X]$,
so $\sum_{j\in J_b}c_jF_{ij}(X)=0$.
The vector equal to $c_j$ on $J_b$ and zero elsewhere is
therefore in $\ker A_i$. Since that kernel is generated by
the full-support vector $c$, every nonempty $J_b$ must be
all of $[w]$. In particular all $b_{\ell j}$ equal
$b_{\ell1}=0$.

It follows that every off-$i$ factor is a constant multiple
of its first-column factor. Equation (I.71) then implies that
$f$ is the product of these common off-$i$ factors with the
zero polynomial $\sum_jc_jF_{ij}(X_i)$, up to a common
nonzero scalar. Thus $f=0$, a contradiction. This proves
the nullity-one assertion. The proof includes the case where
a proposed one-dimensional kernel has a zero coordinate:
then (I.71) cannot occur, and the $i$th cylinder is already
nonzero.

\paragraph{3. Three terms and the general local-nullity reduction.}
For three nonzero product terms, either some $A_i$ has rank
at least two, when the preceding argument applies, or every
$A_i$ has rank one. In the second case every factor list is
proportional, and $f$ is a scalar multiple of one separated
product. If nonzero, it survives the base curve. The cases
of one or two terms follow directly as well. Hence every
nonzero sum of at most three products has a nonzero cylinder
with at most one free coordinate.

For later use, we prove by induction that a sum of at most
$s\ge3$ separated products has a nonzero cylinder with at
most $s-2$ free coordinates (or all $n$ coordinates if fewer
are available). The base $s=3$ was just proved. If every
local factor matrix has rank one, the assertion again
reduces to one product. Otherwise choose a coordinate $i$
whose local rank is $r\ge2$, and choose $r$ of its factors
as a basis. Expanding the other factors in this basis gives
\[
 f=\sum_{a=1}^r F_{i,j_a}(X_i)Q_a(\mathbf X_{\ne i}),
 \tag{I.74}
\]
where each $Q_a$ is a sum of at most $s-r+1\le s-1$
separated products: the original pivot term and at most one
term for each nonpivot column. If every cylinder of $f$
with at most $s-2$ free coordinates were zero, independence
of the basis factors would imply that every $Q_a$ vanished
on cylinders with at most $s-3$ free coordinates. The
induction hypothesis would force every $Q_a=0$, contrary
to $f\ne0$. The prime threshold inherited by a subproblem
on fewer variables is still sufficient.

Now suppose a specified $A_i$ has nullity $c\ge2$.
Expansion (I.74) makes each $Q_a$ a sum of at most $c+1$
products. The result just proved detects each nonzero $Q_a$
with at most $c-1$ free coordinates. Adding $i$ gives at
most $c$ free coordinates for $f$. Nullity zero is handled
by freeing $i$ alone, and nullity one by paragraph~2. This
proves all the stated cylinder bounds and, by paragraph~1,
the explicit cocharacter and sampling bounds.

The same finite family jointly preserves the coefficient
span of the vector of product terms under the local-nullity
promise. More precisely, if $V\subseteq F^w$ is that
coefficient span and the sampling grid lies in a working
extension $E$, the sampled vectors span $V\otimes_F E$
over $E$. Indeed, every $E$-linear combination is still
such a sum over $E$ after scaling or deleting columns of
$A_i$; column deletion cannot increase its nullity, and
scalar extension preserves the original matrix rank.
If an $E$-linear functional annihilates every sampled vector,
the scalar hitting theorem applied over $E$ forces it to
annihilate $V\otimes_F E$. The reverse inclusion follows
because every evaluated vector is a combination of the
original coefficient vectors. In particular no original-field
functional nonzero on $V$ is lost. This is the exact
joint-span criterion I-14. We do not assert that one member
of the family is always a BIWA or preserves the entire span
individually.

\paragraph{4. An actual additive cancellation of three dense products.}
The multiplicative conclusion in I-33 does not extend to an
additive relation among three products. Let $n=6$, $\delta=12$,
and assign the following distinct prime weights to variables
$A,B,C,Q,R,S$, respectively:
\[
 (1663,523,139,127,1093,331).
\]
All exceed $H=120$ and are below the polynomial bound above.
The polynomial
\[
 \begin{aligned}
 F={}&(1-A)(1-S^2)-(1-R^2)(1-C)\\
     &{}-C(1-B)(1-Q^{12})
 \end{aligned}
\]
is nonzero: the monomial $AS^2$ has coefficient one.
Its three displayed terms are separated products, and no
pairwise ratio is a monomial in the original variables.
Nevertheless this prime curve kills it identically in every
characteristic. To verify this, put $a=1663$, $b=523$,
$c=139$. The remaining weights satisfy
$2\cdot331=b+c$, $2\cdot1093=a+b$, and
$12\cdot127=a-c$. The elementary identity
\[
 (1-U)(1-VW)-(1-UV)(1-W)=(1-V)(W-U)
\]
with $U=T^a$, $V=T^b$, $W=T^c$ proves the assertion.

The repaired family detects this example. Freeing $A$ and
holding the other variables on their prime curve gives
$(T^{1663}-A)(1-T^{662})\ne0$. This example refutes the
specific proposed extension of the divisor lemma; it is
not a lower bound for sparse PIT or for the repaired family.

\paragraph{5. Nullity two already defeats every single-coordinate repair.}
Let $n=8$, $\delta=1$, and take the first eight primes
greater than $H=14$:
\[
 (17,19,23,29,31,37,41,43).
\]
Set
\[
 F=(X_1X_8-X_2X_7)(X_3X_6-X_4X_5).
\]
Both binomials vanish on the base curve, because
$17+43=19+41=60$ and $23+37=29+31=60$.
On any cylinder freeing one coordinate, at least one of
the two disjoint binomials remains unchanged and zero.
Thus every one-coordinate cylinder is identically zero,
including all substitutions that change that coordinate
to any alternate prime power or any rational function.

The polynomial has tensor rank exactly four in the separated
product model. Expansion gives the upper bound four. For
the lower bound, partition the variables as
$\{1,8,3,6\}$ and $\{2,7,4,5\}$.
Each binomial has a rank-two coefficient matrix across its
corresponding partition, and the product on disjoint
variables gives their Kronecker product, of rank four.
Every separated product has coefficient-matrix rank one,
so at least four terms are necessary. This argument works
in every characteristic. In the displayed four-term
representation each local factor matrix has rank two and
nullity two.

Freeing, for example, $X_1$ and $X_3$ makes both binomials
nonzero, so the two-coordinate repair proved above removes
this precise obstruction. The next unresolved efficiency
question is to avoid a sampling exponent proportional to
unbounded local nullity. Polynomially many alternate primes
alone cannot do that while only one coordinate changes.
No conversion of this repair into a Hoza path-isolation
construction has been proved.\hfill$\square$

\paragraph{Derivation record.}
The full proof is a deduction in this track from the I-33
divisor lemma and elementary coefficient elimination.
The specific prime numbers, the three-term identity, and
the tensor-rank-four obstruction were checked directly.
No literature-priority claim is made.
